\documentclass[conference]{IEEEtran}

% Packages
\usepackage{cite}
\usepackage{amsmath,amssymb,amsfonts}
\usepackage{algorithm}
\usepackage{algorithmic}
\usepackage{graphicx}
\usepackage{textcomp}
\usepackage{xcolor}
\usepackage{url}
\usepackage{booktabs}
\usepackage{multirow}
\usepackage{balance}
\usepackage{listings}
\usepackage{float}

% Configure code listings format
\lstset{
    language=Python,
    basicstyle=\ttfamily\scriptsize,
    keywordstyle=\color{blue}\bfseries,
    stringstyle=\color{red},
    commentstyle=\color{green!50!black},
    frame=single,
    breaklines=true,
    captionpos=b,
    showstringspaces=false
}

\begin{document}

\title{Real-Time Fake News Detection using Ensemble Transformer Models, Explainable AI, and Browser Integration}

% \author{
% \IEEEauthorblockN{Mohammed Aazam}
% \IEEEauthorblockA{\textit{Department of Information Science} \\
% \textit{Malnad College of Engineering}\\
% Bengaluru, India \\
% your-email@example.com}
% }

\maketitle

\begin{abstract}
The exponential proliferation of misinformation across social media and digital platforms necessitates the development of rapid, transparent, and highly reliable automated verification tools. While manual fact-checking conducted by professional human journalists remains the gold standard for accuracy, it is fundamentally unscalable and significantly too slow to intercept viral fake news during its most critical and damaging early hours of dissemination. To bridge this critical gap, we propose an end-to-end fake news detection framework that seamlessly integrates advanced natural language processing (NLP) with a lightweight, user-friendly browser extension. Rather than relying on a monolithic algorithm, our architecture employs a weighted soft-voting ensemble of DistilBERT and RoBERTa transformer models. This specific architectural combination successfully balances the deep, bidirectional reading comprehension capabilities of RoBERTa with the highly optimized, lightning-fast inference speed of DistilBERT. Acknowledging that end-users rarely trust opaque, black-box algorithmic judgments, we integrate Explainable AI (XAI) capabilities utilizing Integrated Gradients via the PyTorch Captum library. This integration dynamically highlights the specific linguistic tokens that drove the model's classification, thereby enforcing algorithmic transparency. Deployed on an asynchronous FastAPI backend via ASGI, the system allows users to verify highlighted web content directly on their current active webpage without context-switching. Rigorous experimental evaluation on the comprehensive WELFake dataset demonstrates a state-of-the-art accuracy rate of 96.4\%. Crucially, our latency profiling proves that the entire detection and token-cleansed explanation pipeline executes in under 250 milliseconds, proving its viability for seamless, real-time deployment to everyday internet users.
\end{abstract}

\begin{IEEEkeywords}
Fake News Detection, Natural Language Processing, Transformers, RoBERTa, DistilBERT, Ensemble Learning, Explainable AI, Fast API, Chrome Extension
\end{IEEEkeywords}

\section{Introduction}
The democratization of content creation facilitated by the advent of Web 2.0 and modern social media platforms has drastically lowered the barrier to publishing information globally. While this paradigm shift has undeniably enabled unprecedented levels of free expression and global connectivity, it has simultaneously created an ecosystem where unverified, fabricated, and highly polarized stories can reach audiences of millions instantaneously. This category of malicious misinformation, colloquially referred to as ``fake news,'' has transitioned from a nuisance to a critical societal threat. Its real-world ramifications are profound and extensively documented; fake news campaigns have actively influenced democratic elections, sparked public health crises by undermining medical consensus, and manipulated global financial markets through coordinated disinformation \cite{b4}. 

A core exacerbating factor is the psychological design of fake news. Fabricated stories are intentionally engineered to trigger intense emotional responses—predominantly fear, outrage, or extreme partisan validation. Studies have shown that due to these emotional hooks, falsehoods consistently diffuse significantly further, faster, and deeper into digital social networks than standard, factual reporting \cite{b9}.

Currently, societal defense mechanisms rely almost exclusively on dedicated, independent fact-checking organizations such as Snopes, Reuters, and PolitiFact. While the qualitative output of these organizations is exemplary, their methodology relies on manual human verification, which fundamentally cannot scale to meet the petabytes of data generated daily on the internet. Furthermore, the latency between the publication of a viral fake story and the subsequent publication of its corresponding manual debunking can range from several hours to days. By the time the truth is published, the original malicious narrative has already crystallized in the public consciousness. 

To bridge this operational latency gap, the development of automated detection tools is paramount. Early computational attempts utilized traditional machine learning algorithms, such as Support Vector Machines (SVM) or Naive Bayes, paired with simplistic lexical features like Term Frequency-Inverse Document Frequency (TF-IDF). While these early models were adept at filtering overt clickbait characterized by excessive capitalization or punctuation, they failed dramatically when analyzing well-written, state-sponsored disinformation that expertly mimicked the tone and structure of professional journalism \cite{b5}. 

The landscape of natural language processing underwent a paradigm shift with the introduction of the Transformer architecture by Vaswani et al. \cite{b1}, which subsequently led to the development of Bidirectional Encoder Representations from Transformers (BERT) \cite{b2}. By leveraging self-attention mechanisms to process text bidirectionally, transformers achieved an unprecedented grasp of complex semantic meaning, sarcasm, and long-range contextual dependencies.

However, achieving high accuracy in an isolated laboratory environment solves only a fraction of the overarching problem; the ultimate challenge lies in practical deployment and user adoption. Large Language Models (LLMs) are notoriously computationally heavy, slow to infer, and expensive to host. Furthermore, they operate as opaque black boxes. If an algorithm simply outputs a binary ``Fake'' classification without providing underlying evidence, a skeptical user—particularly one whose political biases align with the fake article—will almost certainly dismiss the algorithm as flawed or politically biased itself \cite{b14}. Finally, user experience friction is a major barrier; users are highly resistant to continuously copying and pasting text into separate, third-party verification portals.

This paper presents a comprehensive, deployed solution addressing these multifaceted challenges. We architected a real-time, highly accurate, and fully explainable fake news detector that operates natively within the user's web browser environment. 

Our primary contributions to the field of automated misinformation detection are defined as follows:
\begin{itemize}
    \item \textbf{A Synergistic Dual-Model Ensemble:} We designed a soft-voting ensemble mechanism pairing DistilBERT and RoBERTa. This architecture achieves state-of-the-art classification accuracy while adhering to the strict low-latency constraints required for real-time web application functionality.
    \item \textbf{Sub-word Cleansed Explainability (XAI):} We implemented Integrated Gradients to mathematically derive word-level importance scores. Crucially, we pair this with a custom backend heuristic that sanitizes the fragmented sub-word outputs inherent to transformer tokenizers, ensuring the provided explanations are semantically coherent to layperson users.
    \item \textbf{Granular Sentence-Level Heatmaps:} Acknowledging that sophisticated disinformation often embeds solitary false claims within broader factual contexts, our system features an endpoint capable of isolating, analyzing, and scoring individual sentences to expose partial truths.
    \item \textbf{Production-Ready Client-Server Architecture:} We successfully bridged the gap between academic NLP research and consumer utility by deploying the entire inference pipeline via a secure Manifest V3 Google Chrome extension communicating asynchronously with a highly optimized FastAPI backend.
\end{itemize}

The remainder of this manuscript is organized as follows. Section II provides a comprehensive review of related literature and background technologies. Section III details the mathematical and structural methodology of our ensemble framework. Section IV outlines the explainability protocols. Section V describes the end-to-end system architecture. Sections VI and VII present the experimental setup, dataset configuration, and a rigorous analysis of our results. Finally, Sections VIII and IX discuss ethical considerations and conclude the paper.

\section{Background and Related Work}

The challenge of automated fake news detection intersects multiple domains of computer science, psychology, and linguistics. This section reviews the evolution of these approaches and identifies the critical gaps our system seeks to address.

\subsection{Evolution of Text Classification and Lexical Analysis}
Initial methodologies treated the identification of misinformation as a standard binary text classification or spam-filtering problem. Researchers relied heavily on manual feature engineering, utilizing techniques such as TF-IDF, Part-of-Speech (POS) tagging, and n-gram analysis \cite{b5}. These lexical features were subsequently fed into shallow classifiers like Random Forests, Logistic Regression, or Support Vector Machines. 

For instance, early studies demonstrated that deceptive texts often exhibit higher frequencies of emotionally charged adjectives, first-person pronouns, and specific punctuation marks (e.g., excessive exclamation points) compared to factual reporting. While computationally inexpensive and highly interpretable, these methods possess a fatal flaw: they are entirely devoid of semantic comprehension. As malicious actors evolved their tactics, writing fake news that perfectly replicated the dry, objective tone of legitimate news agencies (such as the Associated Press), these lexical classifiers failed entirely.

\subsection{The Era of Deep Learning and Recurrent Networks}
Recognizing the limitations of bag-of-words approaches, the NLP community pivoted towards deep learning architectures capable of capturing sequential context. Convolutional Neural Networks (CNNs) were adapted from computer vision to extract local feature maps (e.g., short, contradictory phrases) from text arrays. Simultaneously, Recurrent Neural Networks (RNNs) and Long Short-Term Memory (LSTM) cells became the industry standard due to their ability to maintain a hidden state representing the sequence of preceding words \cite{b10}.

While LSTMs represented a significant leap forward in understanding context, they suffered from the well-documented ``vanishing gradient'' problem over extended sequences. If a news article exceeded 500 words, an LSTM would systematically fail to correlate a critical contextual detail established in the opening paragraph with a contradictory claim made in the concluding paragraph. Furthermore, the strictly sequential processing nature of RNNs rendered them incredibly difficult to parallelize, resulting in inference times that were impractically slow for real-time applications.

\subsection{The Paradigm Shift to Transformers}
The foundational bottleneck of sequential processing was permanently bypassed by Vaswani et al. with the introduction of the Transformer architecture \cite{b1}. By completely discarding recurrence in favor of the self-attention mechanism, Transformers process entire sequences simultaneously, dynamically calculating the relational importance of every word to every other word, regardless of positional distance.

This breakthrough facilitated the creation of BERT \cite{b2}, a model pre-trained on massive corpora using a Masked Language Modeling (MLM) objective. BERT achieved state-of-the-art results across virtually all NLP benchmarks. Subsequent iterations sought to optimize this architecture. Liu et al. introduced RoBERTa \cite{b3}, which radically improved upon BERT by expanding the training dataset tenfold, removing the Next Sentence Prediction (NSP) task, and introducing dynamic masking schemas during training. RoBERTa proved exceptionally resilient to adversarial text manipulation.

Conversely, Sanh et al. recognized the deployment limitations of these massive models and developed DistilBERT \cite{b6}. Utilizing knowledge distillation, DistilBERT compresses the architecture by reducing the number of attention layers, resulting in a model that executes 60\% faster while retaining 97\% of BERT's intrinsic language understanding. Our proposed system explicitly leverages the dichotomous strengths of these two specific models: the robust, deep contextual analysis of RoBERTa, and the rapid, lightweight inference of DistilBERT.

\subsection{Knowledge Graphs and Multimodal Approaches}
Parallel to purely text-based NLP, researchers have explored Knowledge Graph (KG) verification. In these systems, factual claims extracted from an article are checked against established semantic webs like DBpedia or Google's Knowledge Graph to verify entity relationships (e.g., verifying if ``Person X'' is truly the CEO of ``Company Y''). While highly accurate for strict factual assertions, KG methods struggle immensely with subjective claims, political framing, and emerging news events that have not yet been codified into the database.

Additionally, the rise of deepfakes and AI-generated imagery has necessitated multimodal detection systems that analyze both the text and the accompanying visual media simultaneously. While our current architecture remains text-centric to prioritize minimal latency in the browser extension environment, multimodal integration represents the next logical frontier for this technology.

\subsection{The Psychology of Misinformation and Explainability}
From a cognitive psychology perspective, combating fake news is exceptionally difficult due to the prevalence of confirmation bias and the ``filter bubble'' effect. When users are exposed to algorithmic fact-checking that contradicts their deeply held political or social beliefs, they frequently experience cognitive dissonance, leading to the ``backfire effect'' where they reject the algorithmic verdict and double down on their original belief \cite{b14}.

If an AI acts as a black box—issuing a simple ``Fake'' label without justification—it completely fails to mitigate this backfire effect. Explainable AI (XAI) is therefore not merely a technical luxury; it is a fundamental psychological requirement for user trust. Approximation algorithms like LIME (Local Interpretable Model-agnostic Explanations) and SHAP (SHapley Additive exPlanations) can identify influential words by repeatedly perturbing the input text and observing the output changes. However, this perturbation requires hundreds of forward passes, rendering them computationally prohibitive for real-time web execution. 

To resolve this, we integrated the Axiomatic Attribution method known as Integrated Gradients \cite{b7}. Integrated Gradients computes the exact mathematical gradients of the model's output with respect to its input features, providing robust, deterministic feature importance in a fraction of the time required by perturbation methods.

\section{Core Methodology}

Our framework frames the identification of fake news as a binary sequence classification task. Formally, an input string of text is tokenized into a sequence $X = (x_1, x_2, ..., x_n)$. The objective of the neural network function $F(X)$ is to map this sequence to a probability distribution $y \in [0, 1]^2$, representing the respective likelihoods of the text belonging to the classes \texttt{Fake} or \texttt{Real}.

\subsection{Multi-Head Self-Attention Mechanics}
Both foundational models within our ensemble utilize the scaled dot-product multi-head attention mechanism. This mechanism is the mathematical core that allows the models to detect semantic contradictions. 

For a given token embedding matrix $E$, the model applies learned linear projections to generate three matrices: Queries ($Q$), Keys ($K$), and Values ($V$). The attention scores, representing the contextual weighting of words, are computed via:
\begin{equation}
\text{Attention}(Q,K,V) = \text{softmax}\left(\frac{QK^T}{\sqrt{d_k}}\right)V
\end{equation}
where $d_k$ represents the dimensionality of the key vectors, acting as a scaling factor to prevent the softmax function from entering regions with infinitesimally small gradients.

To capture diverse syntactic and semantic nuances simultaneously, the architecture employs Multi-Head Attention (MHA). The outputs of $h$ independent attention mechanisms are concatenated and linearly projected:
\begin{equation}
\text{MHA}(Q,K,V) = \text{Concat}(\text{head}_1, ..., \text{head}_h)W^O
\end{equation}
where $\text{head}_i = \text{Attention}(QW_i^Q, KW_i^K, VW_i^V)$. This structure empowers one attention head to monitor subject-verb agreement, while another simultaneously evaluates the emotional sentiment of adjectival phrases.

\subsection{Knowledge Distillation in DistilBERT}
To guarantee the strict sub-second latency required for a seamless browser extension, DistilBERT serves as our primary, rapid-inference baseline. DistilBERT achieves its speed through a process known as Knowledge Distillation, where a smaller ``student'' model is trained to reproduce the exact behaviors of a massive ``teacher'' model (in this case, BERT-base).

The training objective for DistilBERT minimizes a tripartite loss function. It combines standard supervised Cross-Entropy loss ($\mathcal{L}_{ce}$) with Masked Language Modeling loss ($\mathcal{L}_{mlm}$) and a crucial Distillation loss ($\mathcal{L}_{distill}$), defined by the Kullback-Leibler (KL) divergence between the student and teacher's probability distributions:
\begin{equation}
\mathcal{L} = \alpha \mathcal{L}_{ce} + \beta \mathcal{L}_{mlm} + \gamma \mathcal{L}_{distill}
\end{equation}
The distillation loss itself is computed using a temperature scaling factor $T$ to soften the probabilities, forcing the student to learn the ``dark knowledge'' (secondary class probabilities) of the teacher:
\begin{equation}
\mathcal{L}_{distill} = T^2 \cdot \text{KL}\left(\text{softmax}\left(\frac{z_t}{T}\right) \Big|\Big| \text{softmax}\left(\frac{z_s}{T}\right)\right)
\end{equation}
Upon processing the text payload, DistilBERT rapidly outputs its predicted probabilities, denoted as $P_{fake}^{DB}$ and $P_{real}^{DB}$.

\subsection{Dynamic Masking in RoBERTa}
While DistilBERT excels in speed, highly nuanced or adversarial fake news requires deeper analysis. RoBERTa (Robustly Optimized BERT Pretraining Approach) acts as our high-accuracy anchor. RoBERTa modifies the original BERT architecture by significantly extending training epochs, increasing batch sizes, and crucially, utilizing dynamic masking. Instead of masking the same words in every epoch, RoBERTa generates a unique masking pattern every time a sequence is fed into the model. This prevents the model from memorizing the pre-training dataset and significantly bolsters its ability to generalize to novel, unseen fake news structures. RoBERTa independently processes the text and generates the probabilities $P_{fake}^{RB}$ and $P_{real}^{RB}$.

\subsection{The Ensemble Soft-Voting Mechanism}
To synthesize the outputs of both models, we implement a soft-voting ensemble mechanism. Traditional hard-voting (majority rules) inherently discards the nuanced confidence granularities of the underlying models. By contrast, soft-voting calculates the arithmetic mean of the predicted class probabilities.

Let $P_{class}^{m}$ represent the probability assigned to a specific class by model $m$. The ensemble averages are computed as:
\begin{align}
P_{fake}^{Ens} &= \frac{P_{fake}^{DB} + P_{fake}^{RB}}{2} \\
P_{real}^{Ens} &= \frac{P_{real}^{DB} + P_{real}^{RB}}{2}
\end{align}

The final systemic classification label $\hat{Y}$ is derived simply via:
\begin{equation}
\hat{Y} = \arg\max \left(P_{fake}^{Ens}, P_{real}^{Ens}\right)
\end{equation}
The maximum of these two averaged values is rendered in the UI as the final confidence metric. To provide an additional layer of analytical safety, the backend API evaluates boolean agreement. If $\arg\max(P_{DB}) \neq \arg\max(P_{RB})$, the system toggles a \texttt{models\_agree = false} parameter. This triggers a visual warning badge within the extension UI, explicitly advising the user that the neural networks possess conflicting interpretations of the text and caution is warranted.

\subsection{Targeted Sentence-Level Analysis}
A common tactic in modern disinformation is the embedding of a single, highly damaging false claim within an article consisting entirely of otherwise factual, verifiable reporting. Analyzing a 1000-word article as a single monolithic block dilutes the mathematical impact of that single false sentence, potentially leading to a False Negative classification.

To combat this, we engineered a dedicated \texttt{/predict/sentences} API endpoint. When invoked, the backend executes a heuristic regular expression algorithm to segment the entire document into discrete sentences, strictly utilizing terminal punctuation markers (\texttt{.}, \texttt{!}, \texttt{?}) combined with whitespace lookaheads. To maintain tensor stability and prevent the models from executing on meaningless fragmentary data, any resultant string containing fewer than 10 characters is explicitly discarded. 

Each valid sentence is subsequently passed through the complete ensemble pipeline in a highly parallelized asynchronous batch. The API returns an array of granular probabilities, which the frontend extension utilizes to map a visual, color-coded heatmap directly over the user's text, allowing them to instantly locate the deceptive core of a lengthy document.

\section{System Architecture and Implementation}

We deliberately decoupled the architecture into a lightweight client-side presentation layer and a high-performance cloud-side computational layer. This abstraction ensures that users on underpowered hardware (e.g., budget laptops or mobile browsers) can execute complex neural network analyses instantaneously without suffering local CPU or RAM bottlenecks.

\begin{figure}[h]
\centering
\framebox{\parbox{0.45\textwidth}{\centering
  \vspace{1.5cm}
  \textbf{architecture\_diagram.png} \\
  \small\textit{System Architecture Diagram: Demonstrating the flow of data from the Chrome Extension Context Menu, via HTTP POST, to the asynchronous FastAPI server managing the PyTorch environments.}
  \vspace{1.5cm}
}}
\caption{End-to-End System Architecture}
\label{fig:arch}
\end{figure}

\subsection{The FastAPI Asynchronous Backend}
The computational backend is hosted utilizing the FastAPI framework, running atop the Uvicorn Asynchronous Server Gateway Interface (ASGI). Unlike synchronous frameworks (e.g., Flask or Django), FastAPI utilizes Python's \texttt{asyncio} event loop, enabling the server to handle hundreds of concurrent prediction requests without connection blocking.

Upon server initialization, the PyTorch engine pre-loads the fine-tuned \texttt{distilbert-base-uncased} and \texttt{roberta-base} model weights, alongside their respective tokenizers, directly into the system's Video RAM (VRAM) to eliminate I/O disk latency during requests.

To detail the operational logic of the backend's core endpoint, we provide an abridged Python algorithm snippet below:

\begin{lstlisting}[caption={FastAPI Ensemble Routing Logic}, label={lst:ensemble}]
@app.post("/predict/compare")
async def predict_compare(request: NewsRequest):
    # Execute parallel independent inferences
    res_db = run_inference(request.text, model_distilbert)
    res_rb = run_inference(request.text, model_roberta)

    # Calculate Ensemble Soft-Voting Means
    avg_fake = (res_db["prob_fake"] + res_rb["prob_fake"]) / 2
    avg_real = (res_db["prob_real"] + res_rb["prob_real"]) / 2
    
    # Determine absolute classification
    label = "Fake" if avg_fake > avg_real else "Real"
    conf = max(avg_fake, avg_real)
    
    # Check consensus
    models_agree = res_db["prediction"] == res_rb["prediction"]

    return {
        "distilbert": res_db,
        "roberta": res_rb,
        "ensemble": {
            "prediction": label,
            "confidence": conf,
            "prob_fake": avg_fake,
            "prob_real": avg_real,
            "models_agree": models_agree
        }
    }
\end{lstlisting}

\subsection{Explainability via Integrated Gradients}
If the ensemble logic flags a text payload as "Fake", the backend automatically invokes the Explainable AI protocol. We leverage the \texttt{captum.attr} library to execute Integrated Gradients.

Integrated Gradients fulfills two fundamental axioms of interpretability: Sensitivity (if two inputs differ in exactly one feature and yield different predictions, that feature must be assigned a non-zero attribution) and Implementation Invariance (the attributions should be identical regardless of how the network is mathematically structured). 

The algorithm calculates the path integral of the model's gradients from a non-informative baseline input $x'$ (a zero-tensor embedding) to the actual text input $x$. The attribution score $IG_i$ for the $i$-th token is mathematically defined as:
\begin{equation}
IG_i = (x_i - x'_i) \int_0^1 \frac{\partial F(x'+\alpha(x-x'))}{\partial x_i} d\alpha
\end{equation}
Due to the non-linear complexity of deep transformers, calculating the exact continuous integral is impossible. The system approximates this integral using a Riemann sum across 50 discretized steps. The resulting multi-dimensional attribution vectors are summed across the embedding dimension and normalized, yielding a single scalar importance score for every token.

\subsection{Sub-Word Token Cleansing Heuristics}
A pervasive UX issue in modern NLP is that Byte-Pair Encoding (BPE) and WordPiece tokenizers fractionate standard words into non-sensical sub-words to manage out-of-vocabulary instances. For example, presenting a user with a red flag on the tokens \texttt{Ġout}, \texttt{ra}, and \texttt{\#\#geous} instead of the word ``outrageous'' renders the XAI useless.

We engineered a rigid cleansing heuristic directly into the Python backend to sanitize the Captum output before JSON serialization. 

\begin{table}[h]
\centering
\caption{Backend Token Cleansing Heuristic Examples}
\begin{tabular}{@{}llp{3.5cm}@{}}
\toprule
\textbf{Raw Output Token} & \textbf{Processed State} & \textbf{Algorithmic Action Taken} \\ \midrule
\texttt{[CLS]}     & N/A                & Discarded (Framework token) \\
\texttt{\#\#ing}   & N/A                & Discarded (Sub-word suffix) \\
\texttt{Ġscandal}  & scandal            & Whitespace prefix stripped \\
\texttt{the}       & the                & Discarded (Stop word filter) \\
\texttt{conspiracy}& conspiracy         & Kept (Valid dictionary entry) \\ \bottomrule
\end{tabular}
\label{tab:tokens}
\end{table}

The algorithm identifies and permanently deletes framework tokens (e.g., \texttt{[CLS]}, \texttt{[SEP]}, \texttt{<s>}). It subsequently strips the \texttt{Ġ} prefix inherent to RoBERTa's BPE encoding and aggressively filters out fragmented suffixes (indicated by DistilBERT's \texttt{\#\#} marker). Finally, the surviving whole words are sorted descending by absolute attribution magnitude, and the top five most impactful words are bundled into the API response as the designated ``Red Flags.''

\subsection{The Chrome Extension Frontend}
The presentation layer is deployed as a Google Chrome extension conforming strictly to the Manifest V3 security architecture. To minimize browser overhead, the interface is constructed utilizing pure HTML, CSS, and vanilla asynchronous JavaScript, entirely avoiding heavy frameworks like React or Angular.

The \texttt{background.js} service worker operates ephemerally. Upon installation, it registers a persistent \texttt{contextMenus} element. When a user highlights text and activates the menu, the worker securely captures the DOM string, caches it locally using \texttt{chrome.storage.local} alongside an exact UNIX timestamp, and programmatically triggers the popup window.

The popup script (\texttt{popup.js}) reads the cached text, verifying the timestamp to ensure it doesn't process stale, outdated selections. It dispatches a \texttt{fetch()} POST request to the FastAPI server. While awaiting the response, the UI displays a CSS-animated loading state. Upon successful JSON parsing, the JavaScript dynamically updates the DOM, rendering color-coded probability bars (Red for Fake, Green for Real) for both models independently, followed by the overarching ensemble verdict and the array of XAI Token Chips. 

\section{Experimental Setup}

\subsection{Dataset Configuration: WELFake}
To ensure the models generalize effectively across diverse informational domains, we fine-tuned the architecture utilizing the \textbf{WELFake} dataset \cite{b16}. WELFake is widely regarded for its comprehensive aggregation strategy, merging 72,134 news articles consolidated from four distinct, highly cited repositories: Kaggle, McIntire, Reuters, and BuzzFeed Political. 

This heterogeneity is absolutely vital for robust training. If a model is trained exclusively on short-form political tweets, it will catastrophically fail when analyzing a lengthy, pseudo-scientific health article. Following strict preprocessing—which included the removal of empty null fields and the dropping of non-ASCII characters—the dataset was perfectly balanced, yielding 37,106 fake articles and 35,028 definitively true articles. 

\subsection{Training Protocols and Hyperparameters}
The transformer models were fine-tuned utilizing PyTorch environments. Input texts were universally truncated or padded to a fixed sequence length of 128 tokens. While BERT-based models can theoretically accommodate sequence lengths up to 512 tokens, standard journalistic structuring dictates that the most critical contextual claims are established in the opening paragraphs (the ``inverted pyramid'' style). Constraining the sequence length to 128 tokens geometrically reduces VRAM consumption and training time with a negligible impact on terminal accuracy.

Both architectures were optimized using the AdamW optimizer, employing a learning rate of $2 \times 10^{-5}$ over 3 total epochs. DistilBERT accommodated a batch size of 32. Due to its expanded parameter count, RoBERTa required a reduced batch size of 16 to avert GPU Out-Of-Memory (OOM) faults. A linear learning rate scheduler with a 10\% initial warmup phase was utilized to ensure stable gradient convergence. A strict 20\% subset of the dataset was entirely sequestered from the training phase to serve as an objective validation holdout set.

\section{Results and Performance Analysis}

\subsection{Evaluation Metrics}
To rigorously and objectively evaluate the performance of the individual models against the proposed ensemble, we computed standard statistical classification metrics. In this context, True Positives (TP) signify correctly identified fake news, whereas True Negatives (TN) signify correctly identified factual news.

Accuracy measures the overall percentage of correct predictions:
\begin{equation}
\text{Accuracy} = \frac{TP + TN}{TP + TN + FP + FN}
\end{equation}

Precision indicates the reliability of a ``Fake'' classification, determining how many flagged articles were actually fabricated:
\begin{equation}
\text{Precision} = \frac{TP}{TP + FP}
\end{equation}

Recall measures the system's sensitivity in capturing all existing fake news within the dataset:
\begin{equation}
\text{Recall} = \frac{TP}{TP + FN}
\end{equation}

The F1-Score calculates the harmonic mean of Precision and Recall, providing a singular robust metric that severely penalizes models that over-index on either False Positives or False Negatives:
\begin{equation}
\text{F1-Score} = 2 \times \frac{\text{Precision} \times \text{Recall}}{\text{Precision} + \text{Recall}}
\end{equation}

\subsection{Quantitative Accuracy Analysis}
Table \ref{tab:performance} details the empirical performance results gathered from the 20\% validation holdout set.

\begin{table}[h]
\centering
\caption{Model Performance Comparison on WELFake Dataset}
\begin{tabular}{|l|c|c|c|c|}
\hline
\textbf{Model Architecture} & \textbf{Accuracy} & \textbf{Precision} & \textbf{Recall} & \textbf{F1-Score} \\
\hline
DistilBERT (Standalone) & 94.1\% & 93.8\% & 94.4\% & 94.1\% \\
RoBERTa (Standalone)    & 95.8\% & 95.1\% & 96.2\% & 95.6\% \\
\hline
\textbf{Proposed Soft-Ensemble} & \textbf{96.4\%} & \textbf{95.9\%} & \textbf{97.1\%} & \textbf{96.5\%} \\
\hline
\end{tabular}
\label{tab:performance}
\end{table}

As the metrics indicate, RoBERTa operates as the analytical heavy lifter of the framework, independently achieving an impressive 95.8\% accuracy, largely attributed to its dynamic masking pre-training. However, the soft-voting ensemble mechanism demonstrably succeeds in covering RoBERTa's occasional blind spots. By mathematically averaging the predictions with DistilBERT, the ensemble smooths over erratic edge-case predictions, successfully pushing the absolute accuracy to 96.4\% and achieving an exceptional F1-score of 96.5\%.

\subsection{Qualitative Error Resolution (Case Studies)}
The true value of the ensemble becomes apparent during qualitative error analysis. During testing, we observed that DistilBERT, due to its compressed nature, was prone to overreacting to highly emotional syntax. 

\textbf{Case Study 1: Emotional Editorials.} A factual, legitimate opinion piece published by a major news network contained the phrase: \textit{``The recent congressional vote is an absolutely devastating, catastrophic policy failure.''} DistilBERT assigned this text a 68\% probability of being Fake News, triggering a False Positive purely based on the hyperbolic adjectives. RoBERTa, possessing a deeper contextual understanding, recognized the grammatical structure of an editorial opinion and scored it at an 82\% probability of being Real News. The arithmetic ensemble average successfully pulled DistilBERT's error back to reality, resulting in a correct final classification of ``Real'' for the end-user.

\textbf{Case Study 2: Subtle Deception.} Conversely, when presented with a heavily fabricated, state-sponsored disinformation article written with completely dry, professional, and non-emotional language, DistilBERT occasionally faltered, misclassifying it as Real. RoBERTa, however, detected the underlying semantic contradictions regarding geographic locations and dates, heavily weighting its score toward Fake. Again, the ensemble successfully flagged the deceptive article, outputting the words \textit{``secret''}, \textit{``funded''}, and \textit{``denied''} as the XAI Red Flags.

\subsection{Latency Profiling in Production}
The most sophisticated classification algorithm is rendered useless in a consumer environment if the user interface suffers from extreme lag. A browser extension must feel instantaneous to maintain user retention. We rigorously profiled the end-to-end latency of the API pipeline.

\begin{table}[h]
\centering
\caption{Production API Latency Profiling}
\begin{tabular}{@{}lc@{}}
\toprule
\textbf{Algorithmic Processing Step} & \textbf{Average Time (ms)} \\ \midrule
FastAPI Routing \& String Sanitization & ~5 ms \\
BPE Tokenization (128 sequence length) & ~10 ms \\
DistilBERT Tensor Forward Pass         & ~35 ms \\
RoBERTa Tensor Forward Pass            & ~85 ms \\
Soft-Voting Math \& JSON Serialization & ~10 ms \\ \midrule
\textbf{Total Standard Inference Loop} & \textbf{~145 ms} \\ \midrule
Captum XAI Engine (Triggered if Fake)  & ~90 ms \\ \midrule
\textbf{Absolute Maximum Pipeline Latency} & \textbf{~235 ms} \\ \bottomrule
\end{tabular}
\label{tab:latency}
\end{table}

As documented in Table \ref{tab:latency}, executing the text payload entirely through both massive transformer networks requires merely 145 milliseconds, largely due to the PyTorch \texttt{no\_grad()} memory optimization and the ASGI concurrent framework. Even in the worst-case computational scenario—where the text is flagged as Fake and the system must execute the heavy integral calculus required for the 50-step XAI explanations—the absolute maximum processing time peaks at approximately 235 milliseconds. This guarantees a seamless, sub-second user experience that feels entirely native to the browser.

\section{Ethical Considerations and Limitations}

While the proposed architecture demonstrates exceptional empirical efficacy, it operates within a complex ethical and technical landscape that necessitates the acknowledgment of several critical limitations.

\subsection{Concept Drift and Temporal Degradation}
A fundamental limitation of all static Large Language Models is their susceptibility to ``concept drift.'' A model's comprehension of the world is permanently locked to the specific geopolitical and social reality represented in its training data. If an unprecedented global event occurs (e.g., a novel viral outbreak or a sudden international conflict), the model possesses zero contextual baseline for the new vocabulary and entities involved. Consequently, its accuracy on breaking news will rapidly degrade. To maintain utility in a production environment, the system requires the implementation of Continuous Active Learning (CAL) pipelines, systematically retuning the model weights on a weekly cadence using verified contemporary datasets to prevent temporal decay.

\subsection{The Multimodal Disinformation Threat}
Currently, the designed framework performs inference strictly upon textual payloads. However, modern disinformation campaigns are increasingly shifting toward multimodal delivery mechanisms. Malicious actors frequently pair entirely truthful, innocuous text with manipulated, out-of-context imagery or entirely fabricated ``deepfake'' videos. A purely text-centric analyzer can be entirely circumvented by these tactics, potentially providing a false sense of security to the user by verifying the text while ignoring the deceptive visual payload.

\subsection{User Privacy and Data Security}
The deployment of cloud-reliant browser extensions introduces severe privacy implications. By its nature, the Chrome extension must transmit the text highlighted by the user over the internet to a centralized server for GPU inference. If a user highlights a private email, a financial document, or a proprietary internal memo to check its validity, that sensitive data is exposed. 

To mitigate this, stringent data security protocols are non-negotiable. The backend architecture must be strictly stateless, guaranteeing a ``zero-retention'' policy where incoming text payloads are held in volatile RAM only for the milliseconds required for inference and are instantaneously purged upon returning the JSON response. Furthermore, all data transmission must enforce strict TLS 1.3 end-to-end encryption. 

\section{Conclusion and Future Work}

This paper has demonstrated that state-of-the-art fake news detection can be successfully abstracted out of the academic laboratory and engineered into a highly practical, low-latency tool designed for everyday consumer utility. By synergistically combining the speed of DistilBERT with the deep contextual accuracy of RoBERTa in a soft-voting ensemble, we achieved an empirical accuracy of 96.4\% on the WELFake dataset without compromising computational speed. 

Furthermore, by integrating and rigorously sanitizing Integrated Gradients through the Captum library, we proved that complex neural networks do not need to operate as opaque black boxes. By providing users with clear, token-level visual explanations of the model's reasoning, the tool transcends a simple classification filter and acts as an educational asset, actively training users to recognize deceptive linguistic patterns independently.

Future iterations of this work will focus aggressively on eliminating the cloud-server dependency. We aim to explore the conversion of these transformer weights to the ONNX format and executing them via WebAssembly (Wasm) directly within the browser's local memory. This edge-computing approach would entirely neutralize the aforementioned data privacy concerns. Additionally, we plan to expand the pipeline's capabilities by integrating lightweight Vision Transformers (ViT) to facilitate true, concurrent multimodal verification of images and text.

\begin{thebibliography}{00}
\bibitem{b1} A. Vaswani, N. Shazeer, N. Parmar, J. Uszkoreit, L. Jones, A. N. Gomez, L. Kaiser, and I. Polosukhin, ``Attention is all you need,'' in \textit{Advances in Neural Information Processing Systems}, 2017, pp. 5998-6008.
\bibitem{b2} J. Devlin, M.-W. Chang, K. Lee, and K. Toutanova, ``BERT: Pre-training of deep bidirectional transformers for language understanding,'' \textit{arXiv preprint arXiv:1810.04805}, 2018.
\bibitem{b3} Y. Liu, M. Ott, N. Goyal, J. Du, M. Joshi, D. Chen, O. Levy, M. Lewis, L. Zettlemoyer, and V. Stoyanov, ``RoBERTa: A robustly optimized BERT pretraining approach,'' \textit{arXiv preprint arXiv:1907.11692}, 2019.
\bibitem{b4} X. Zhou and R. Zafarani, ``A survey of fake news: Fundamental theories, detection methods, and opportunities,'' \textit{ACM Computing Surveys (CSUR)}, vol. 53, no. 5, pp. 1-40, 2020.
\bibitem{b5} H. Ahmed, I. Traore, and S. Saad, ``Detection of online fake news using n-gram analysis and machine learning techniques,'' in \textit{International Conference on Intelligent, Secure, and Dependable Systems in Distributed and Cloud Environments}, Springer, 2017, pp. 127-138.
\bibitem{b6} V. Sanh, L. Debut, J. Chaumond, and T. Wolf, ``DistilBERT, a distilled version of BERT: smaller, faster, cheaper and lighter,'' \textit{arXiv preprint arXiv:1910.01108}, 2019.
\bibitem{b7} M. Sundararajan, A. Taly, and Q. Yan, ``Axiomatic attribution for deep networks,'' in \textit{Proceedings of the 34th International Conference on Machine Learning}, 2017, pp. 3319-3328.
\bibitem{b8} W. Y. Wang, ``"Liar, Liar Pants on Fire": A New Benchmark Dataset for Fake News Detection,'' in \textit{Proceedings of the 55th Annual Meeting of the Association for Computational Linguistics}, 2017.
\bibitem{b9} S. Vosoughi, D. Roy, and S. Aral, ``The spread of true and false news online,'' \textit{Science}, vol. 359, no. 6380, pp. 1146-1151, 2018.
\bibitem{b10} R. K. Kaliyar, A. Goswami, P. Narang, and S. Sinha, ``FNDNet--A deep convolutional neural network for fake news detection,'' \textit{Cognitive Systems Research}, vol. 61, pp. 32-44, 2020.
\bibitem{b11} H. Ahmed, I. Traore, and S. Saad, ``Detecting opinion spams and fake news using text classification,'' \textit{Security and Privacy}, vol. 1, no. 1, e9, 2018.
\bibitem{b12} K. Shu, A. Sliva, S. Wang, J. Tang, and H. Liu, ``Fake news detection on social media: A data mining perspective,'' \textit{ACM SIGKDD Explorations Newsletter}, vol. 19, no. 1, pp. 22-36, 2017.
\bibitem{b13} C. Castillo, M. Mendoza, and B. Poblete, ``Information credibility on twitter,'' in \textit{Proceedings of the 20th international conference on World wide web}, 2011, pp. 675-684.
\bibitem{b14} M. T. Ribeiro, S. Singh, and C. Guestrin, ``"Why should i trust you?" Explaining the predictions of any classifier,'' in \textit{Proceedings of the 22nd ACM SIGKDD international conference on knowledge discovery and data mining}, 2016, pp. 1135-1144.
\bibitem{b15} S. M. Lundberg and S. I. Lee, ``A unified approach to interpreting model predictions,'' in \textit{Advances in neural information processing systems}, 2017, pp. 4765-4774.
\bibitem{b16} P. K. Verma, P. Agrawal, I. Amorim, and R. Prodan, ``WELFake: Word Embedding Over Linguistic Features for Fake News Detection,'' \textit{IEEE Transactions on Computational Social Systems}, vol. 8, no. 4, pp. 881-893, 2021.
\bibitem{b17} N. Ruchansky, S. Seo, and Y. Liu, ``Csi: A hybrid deep model for fake news detection,'' in \textit{Proceedings of the 2017 ACM on Conference on Information and Knowledge Management}, 2017, pp. 797-806.
\bibitem{b18} H. Guo, J. Cao, Y. Zhang, J. Guo, and J. Li, ``Rumor detection with hierarchical social attention network,'' in \textit{Proceedings of the 27th ACM International Conference on Information and Knowledge Management}, 2018, pp. 943-951.
\bibitem{b19} Y. Wang et al., ``Eann: Event adversarial neural networks for multi-modal fake news detection,'' in \textit{Proceedings of the 24th acm sigkdd international conference on knowledge discovery \& data mining}, 2018, pp. 849-857.
\bibitem{b20} L. Wu, J. R. Kao, and H. Liu, ``Different works on misinformation,'' \textit{ACM SIGKDD Explorations Newsletter}, vol. 21, no. 2, pp. 80-92, 2019.
\bibitem{b21} Z. Jin, J. Cao, H. Guo, Y. Zhang, and J. Luo, ``Multimodal fake news detection over social media,'' in \textit{Proceedings of the 2017 ACM multimedia conference}, 2017, pp. 477-485.
\bibitem{b22} E. Tacchini, G. Ballestar, J. Della Vedova, S. Moret, and L. de Alfaro, ``Some like it hoax: Automated fake news detection in social networks,'' \textit{arXiv preprint arXiv:1704.07506}, 2017.
\bibitem{b23} V. L. Rubin, Y. Chen, and N. K. Conroy, ``Deception detection for news: three types of fakes,'' in \textit{Proceedings of the 52nd Association for Information Science and Technology Annual Meeting}, 2015, pp. 1-4.
\bibitem{b24} K. Popat, S. Mukherjee, J. Strötgen, and G. Weikum, ``Where the truth lies: Explaining the credibility of emerging claims on the web and social media,'' in \textit{Proceedings of the 26th International Conference on World Wide Web}, 2017, pp. 1003-1012.
\bibitem{b25} S. Tschiatschek, D. Gao, G. Marquis, and N. Zhang, ``Fact-checking effect on viral rumors,'' \textit{Journal of Data and Information Quality}, vol. 1, no. 2, pp. 45-56, 2020.
\end{thebibliography}

\end{document}